\subsection{SOME FORMULAS}

In this number, A denotes a unital associative algebra. If \(x \in A\) , we write ad x instead of \(ad_{A}x\) .

Lemma 2. If \(x, y \in \mathbf{A}\) and \(n \in \mathbf{N}\) ,

\[
\frac {(\operatorname{ad} x) ^ {n}}{n !} y = \sum_ {p + q = n} (- 1) ^ {q} \frac {x ^ {p}}{p !} y \frac {x ^ {q}}{q !} = \sum_ {p + q = n} (- 1) ^ {q} x ^ {(p)} y x ^ {(q)}.\tag{11}
\]

Indeed, if we denote by \(L_{x}\) and \(R_{x}\) the maps \(z \mapsto xz\) and \(z \mapsto zx\) from A to A, we have, since \(L_{x}\) and \(R_{x}\) commute,

\[
\frac {1}{n !} (\mathrm{ad} x) ^ {n} = \frac {1}{n !} (\mathrm{L} _ {x} - \mathrm{R} _ {x}) ^ {n} = \sum_ {p + q = n} (- 1) ^ {q} \frac {1}{p !} \mathrm{L} _ {x} ^ {p} \frac {1}{q !} \mathrm{R} _ {x} ^ {q}.
\]

Lemma 3. Let \(x, h \in \mathbf{A}\) and \(\lambda \in k\) be such that (ad \(h)x = \lambda x\) . For all \(n \in \mathbf{N}\) , and all \(\mathrm{P} \in k[\mathrm{X}]\) , we have

\[
\mathrm{P} (h) x ^ {(n)} = x ^ {(n)} \mathrm{P} (h + n \lambda).\tag{12}
\]

Since \(\operatorname{ad} h\) is a derivation of \(A\) and since \((\operatorname{ad} h)x\) commutes with \(x\) , we have

\[
(\operatorname{ad} h) x ^ {n} = n x ^ {n - 1} ((\operatorname{ad} h) x) = n \lambda x ^ {n},\tag{13}
\]

SO

\[
(\operatorname{ad} h) x ^ {(n)} = n \lambda x ^ {(n)}.
\]

Thus, formula (12) follows from the special case

\[
\mathrm{P} (h) x = x \mathrm{P} (h + \lambda)\tag{14}
\]

by replacing \(x\) by \(x^{(n)}\) and \(\lambda\) by \(n\lambda\) . It suffices to prove (14) when \(\mathrm{P} = \mathrm{X}^m\) , by induction on \(m\) . It is clear when \(m = 0, 1\) . If (14) is true for \(\mathrm{P} = \mathrm{X}^m\) , then

\[
h ^ {m + 1} x = h. h ^ {m} x = h x (h + \lambda) ^ {m} = x (h + \lambda) ^ {m + 1}
\]

which proves (12).

Lemma 4. Let \(x, y, h \in \mathbf{A}\) be such that

\[
[ y, x ] = h, \quad [ h, x ] = 2 x, \quad [ h, y ] = - 2 y.\tag{15}
\]

\begin{enumerate}
\def\labelenumi{(\roman{enumi})}
\tightlist
\item
  For \(m, n \in \mathbf{N}\) , we have
\end{enumerate}

\[
x ^ {(n)} y ^ {(m)} = \sum_ {p \geq 0} y ^ {(m - p)} \binom{m + n - p - 1 - h}{p} x ^ {(n - p)}.\tag{16}
\]

\begin{enumerate}
\def\labelenumi{(\roman{enumi})}
\setcounter{enumi}{1}
\tightlist
\item
  Let \(\mathbf{A}'\) be the \(\mathbf{Z}\) -subalgebra of \(\mathbf{A}\) generated by the \(x^{(m)}\) and the \(y^{(m)}\) for \(m \in \mathbf{N}\) . Then \(\binom{h}{n} \in \mathbf{A}'\) for all \(n \in \mathbf{N}\) .
\end{enumerate}

Formula (16) can be written in the equivalent form

\[
(\operatorname{ad} x ^ {(n)}) y ^ {(m)} = \sum_ {p \geq 1} y ^ {(m - p)} \binom {m + n - p - 1 - h} {p} x ^ {(n - p)}.\tag{$17_{m}$}
\]

This is trivial for \(m = 0\) . We argue by induction on \(m\) . From \((17_{m})\) , we obtain

\[
\begin{array}{l} (m + 1) (\operatorname{ad} x ^ {(n)}) y ^ {(m + 1)} = (\operatorname{ad} x ^ {(n)}) y ^ {(m)}. y + y ^ {(m)}. (\operatorname{ad} x ^ {(n)}) y \\ = \sum_ {p \geq 1} y ^ {(m - p)} \binom {m + n - p - 1 - h} {p} x ^ {(n - p)} y + y ^ {(m)} (n - 1 - h) x ^ {(n - 1)} \end{array} \tag {18}
\]

(§1, no. 1, Lemma 1). Now, applying the same lemma, and then Lemma 3, we have

\[
\begin{array}{l} \binom {m + n - p - 1 - h} {p} x ^ {(n - p)} y \\ = \binom {m + n - p - 1 - h} {p} (y x ^ {(n - p)} + (n - p - 1 - h) x ^ {(n - p - 1)}) \\ = y \binom {m + n - p + 1 - h} {p} x ^ {(n - p)} \\ \qquad + \binom {m + n - p - 1 - h} {p} (n - p - 1 - h) x ^ {(n - p - 1)}. \end{array}
\]

Inserting this into (18), we obtain

\[
\begin{array}{l} (m + 1) (\operatorname{ad} x ^ {(n)}) y ^ {(m + 1)} \\ \qquad = \sum_ {p \geq 1} (m - p + 1) y ^ {(m - p + 1)} \binom {m + n - p + 1 - h} {p} x ^ {(n - p)} \\ \qquad + \sum_ {p \geq 1} y ^ {(m - p)} \binom {m + n - p - 1 - h} {p} (n - p - 1 - h) x ^ {(n - p - 1)} \\ \qquad + y ^ {(m)} (n - 1 - h) x ^ {(n - 1)} \\ \qquad = \sum_ {p \geq 1} (m - p + 1) y ^ {(m - p + 1)} \binom {m + n - p + 1 - h} {p} x ^ {(n - p)} \\ \qquad + \sum_ {p \geq 0} y ^ {(m - p)} \binom {m + n - p - 1 - h} {p} (n - p - 1 - h) x ^ {(n - p - 1)}. \end{array}
\]

Changing \(p\) to \(p - 1\) in the second sum, and regrouping the terms, we obtain

\[
(m + 1) (\operatorname{ad} x ^ {(n)}) y ^ {(m + 1)} = \sum_ {p \geq 1} y ^ {(m - p + 1)} A _ {p} x ^ {(n - p)}\tag{19}
\]

with

\[
A _ {p} = (m - p + 1) \binom{m + n - p + 1 - h}{p} + (n - p - h) \binom{m + n - p - h}{p - 1}.
\]

Putting \(z = m + n - p - h\) , this can also be written as

\[
\begin{array}{l} A _ {p} = \frac {1}{p !} (m - p + 1) (z + 1) z (z - 1) \dots (z - p + 2) \\ \qquad + \frac {1}{(p - 1) !} (z - m) z (z - 1) \dots (z - p + 2) \\ \qquad = \frac {1}{p !} z (z - 1) \dots (z - p + 2) [ (m - p + 1) (z + 1) + p (z - m) ] \\ \qquad = (m + 1) \binom {z} {p} = (m + 1) \binom {(m + 1) + n - p - 1 - h} {p}. \end{array}
\]

Inserting this into (19), we obtain \((17_{m + 1})\) , hence (i).

Assume that \(\binom{h}{p} \in \mathrm{A}'\) for \(p < n\) . Then, for all \(\mathrm{P} \in \mathbf{Q}[\mathrm{T}]\) of degree \(< n\) such that \(\mathrm{P}(\mathbf{Z}) \subset \mathbf{Z}\) , we have \(\mathrm{P}(h) \in \mathrm{A}'\) (no. 4, Cor. of Prop. 2). Hence, in view of (16) with \(m = n\) ,

\[
\begin{array}{c} (- 1) ^ {n} \binom{h}{n} = \binom{n - 1 - h}{n} \\ = - x ^ {(n)} y ^ {(n)} + \sum_ {p = 0} ^ {n - 1} y ^ {(n - p)} \binom{2 n - p - 1 - h}{p} x ^ {(n - p)} \in \mathrm{A} ^ {\prime}; \end{array}
\]

hence (ii) by induction on \(n\) .

